% Transcribed from
% https://books.google.com/books?id=8ipKAQAAMAAJ&pg=PA1

\def\itquoted ... „{\it xyz}“
\def\section ... center and letterspace

Avellanedas DON QUIJOTE,
sein Verhältnis zu Cervantes
und seine Bearbeitung durch Lesage.

Von
MARTIN WOLF.

Erster Teil.

Avellaneda und Cervantes.

Nachdem Cervantes neun Jahre lang auf den versprochenen zweiten
Teil seines {\it Don Quijote} hatte warten lassen, erschien im Jahre 1614
ein apokryphischer zweiter Teil unter dem Titel: {\it Segundo Tomo del
ingenioso Hidalgo Don Quixote de la Mancha por el Licenciado Alonso
Fernandez de Avellaneda}. Sicher hätte man diesem Buche in der
Folgezeit ebensowenig Beachtung geschenkt wie anderen minderwertigen
Fortsetzungen von bedeutenden Werken, wenn nicht Cervantes selbst
in seinem zweiten Teil die Nachwelt darauf aufmerksam gemacht hätte.
Das Verbrechen des Plagiats ist jedoch für jene Zeit nicht so selten
und in unserem Falle nicht so gross, wie es nach dem Verhalten des
Cervantes und den Aussprüchen späterer Beurteiler Avellanedas scheinen
möchte. Wir sehen in dem falschen {\it Don Quijote} nur ein neues
Beispiel jener Fortsetzungsmanie, die in der spanischen Litteratur des
16.~und 17.~Jahrhunderts auftritt, und der wenig Werke von einiger
Bedeutung entgangen sind. Die {\it Celestina} des Fernando de Rojas
eröffnet mit etwa zwanzig Fortsetzungen und Nachahmungen den Reigen.
Ihr schliessen sich an der {\it Amadis de Gaula} und andere berühmte
Ritterromane, die {\it Diana} von Montemayor und der {\it Lazarillo de Tormes}.
Unter ganz ähnlichen Umständen wie der falsche {\it Don Quijote} ist eine
Fortsetzung von Mateo Alemans {\it Guzman de Alfarache} entstanden,
die~1602\footnote{1}{Nach Perez Pastor; vgl.\ dazu Morel-Fatio,
{\it Bull. Hisp.}~V,~4, S.~361.}) in Valencia unter dem Pseudonym
Mateo Luxan de Sayavedra
(d.~i.\ Iuan Marti, ein valencianischer Advokat) erschienen ist. Während
sich hier der Entlarvung des Pseudonyms infolge der deutlichen Hinweise
Alemans keine Hindernisse entgegensetzten, ist die Frage nach dem
Verfasser des Pseudo-Don Quijote bis zum heutigen Tage nicht in
befriedigender Weise beantwortet worden. Denn wir haben es auch hier
mit einem Decknamen zu tun. Da das Buch selbst keine Anhaltspunkte
bot, haben die Nachforschungen in der Hauptsache an die Vorrede
Avellanedas und an Cervantes' Äusserungen angeknüpft. So ist es auch
gekommen, dass man das Werk Avellanedas bisher noch nicht auf
seinen litterarischen Wert oder wenigstens auf seine litterarische Stellung
geprüft hat. Dies zu tun, habe ich in dieser Arbeit unternommen und
ich hoffe, dass eine Vergleichung des falschen {\it Don Quijote} mit seinem
Vorbilde auch fruchtbar für die Beurteilung des Originalwerkes sein
wird. Für die französische Litteratur hat Avellanedas Buch einige
Bedeutung, weil Lesage mit einer freien Bearbeitung desselben seine
Laufbahn als Romanschriftsteller eröffnete. Über diesen Erstlingsroman soll
der zweite Teil dieser Arbeit handeln.

\section{1. Aufnahme und Verbreitung.}

Avellanedas Fortsetzung des {\it Don Quijote} scheint keinen Erfolg
gehabt zu haben. Die Zeitgenossen schweigen vollkommen darüber;
selbst in dem umfangreichen Briefwechsel Lope de~Vegas\footnote{2}{Das
Wichtigste ist veröffentlicht in der Biographie von Cayetano Alberto
de la Barrera ({\it Obras} de Lope de~Vega publicadas por la Real Academia Española,
Madrid~1890~ff., I.~Bd.).}) mit dem
Herzog von Sessa, in dem sonst alle litterarischen Ereignisse von einiger
Wichtigkeit ihren Widerhall fanden, wird sie nirgends erwähnt. Das
einzige Zeugnis über Avellanedas Buch aus dem 17.~Jahrhundert ist eine
kurze Notiz in Nicolas Antonio, {\it Bibliotheca Hispana Nova}
(1672--79)\footnote{3}{Zweite Auflage von Francisco Perez Bayer,
Madrid~1783--88: {\itquoted Alphonsus Fernandez de Avellaneda,
patria ex oppido Tordesillas Pincianae diœcesis, continuavit
sed absque genio illo qui principem Michaelis Cervantes ad inventionem
promovit et comitatus est.}}),
die nicht viel mehr enthält, als was schon der Titel des Buches sagt.
Im 17.~Jahrhundert ist keine neue Auflage erschienen, so dass das Buch zur



